\begin{textsample}{POS Dim 4 – 2015\_c6 – Score 12.00 – 2015\_c6/enade\_2015\_tec\_des\_inter\_2e210f18ec.txt}  \label{ex:f4_pos_028}
CAPÍTULO 3 ANÁLISE TÉCNICA DA PROVA Este capítulo tem por objetivo apresentar o desempenho dos estudantes concluintes de Tecnologia em Design de Interiores no Enade/2015.
Para isso, foram calculadas as estatísticas básicas da prova em seu todo ( seção 3.1.1 ), bem como as estatísticas relacionadas aos Componentes de Formação Geral ( seção 3.1.2 ) e de Conhecimento Específico da Área ( seção 3.1.3 ).
Dadas as suas características, foram analisadas, em separado, as questões objetivas ( seção 3.2 ) e as questões discursivas ( seção 3.3 ).
Para as questões objetivas, foram disponibilizados os índices de facilidade e de discriminação bisserial, também em separado, para os Componentes de Formação Geral ( seção 3.2.1 ) e de Conhecimento Específico ( seção 3.2.2 ).
De cada componente, uma das questões foi escolhida para exemplificar a análise gráfica, relacionando as alternativas escolhidas pelos estudantes ( inclusive o gabarito ) com o número de acertos no componente.
O Anexo I apresenta a íntegra da análise gráfica para todas as questões objetivas.
Para cada uma das questões discursivas, os conteúdos dos tipos mais comuns de respostas dos estudantes são apresentados e comparados com o padrão de respostas esperado ( ver Anexo VIII com o padrão de respostas ).
Tomando como base as duas questões discursivas do Componente de Formação Geral, a seção 3.3.1.6 apresenta comentários sobre a correção das respostas com respeito à Língua Portuguesa.
Nas Tabelas 3.1 e 3.3 são apresentados o tamanho da população inscrita, ausente, presente e percentual de faltosos, respectivamente, por Grande Região e por Categoria Administrativa e Organização Acadêmica da IES.
As demais tabelas apresentam as seguintes estatísticas das notas¹⁵ : média do desempenho na prova, erro padrão da média, desvio padrão, nota mínima, mediana e nota máxima para cada um de seus componentes.
Tais estatísticas contemplam o total de estudantes concluintes da Área de Tecnologia em Design de Interiores inscritos e presentes à prova do Enade/2015, tendo em vista agregações ou por Grandes Regiões e o país como um todo, ou por Categoria Administrativa e Organização Acadêmica da IES.
Em relação aos gráficos de distribuição de notas, o intervalo considerado foi de 10 unidades, aberto à esquerda e fechado à direita, com exceção do primeiro intervalo, [ 0 ; 10 ], fechado em ambos os extremos.
Para os gráficos de distribuição das notas das questões discursivas, ¹⁵ Uma definição dessas estatísticas pode ser encontrada no Glossário. questões discursivas, foram consideradas mais duas categorias : questão em branco16 e nota zero. 3.1 CONDIÇÃO DE PRESENÇA E ESTATÍSTICAS BÁSICAS DA PROVA Esta seção apresenta a condição de presença, além das estatísticas selecionadas e histogramas da nota geral ( 3.1.1 ) e de cada componente : Formação Geral ( 3.1.2 ) e Conhecimento Específico ( 3.1.3 ).
São também apresentadas estatísticas selecionadas de subpopulações caracterizadas por Grande Região, Categoria Administrativa e Organização Acadêmica. 3.1.1 Condição de Presença e Estatísticas Básicas Gerais A Tabela 3.1 apresenta a Condição de Presença na prova, por Grande Região, dos estudantes concluintes de Tecnologia em Design de Interiores.
A população total de Inscritos foi de 3.068.
Destes, 2.363 estiveram Presentes, sendo 23,0 \% o índice de não comparecimento.
A Região de maior abstenção foi a Centro-Oeste ( 33,5 \%, ) e a de menor abstenção foi a Sul ( 12,6 \% ).
Já a Tabela 3.2 apresenta as Estatísticas Básicas da prova, por Grande Região, dos estudantes concluintes de Tecnologia em Design de Interiores.
A Média das notas da prova como um todo ( nas seções seguintes serão analisados os Componentes de Formação Geral e de Conhecimento Específico ) foi 46,2, sendo que os alunos da região Norte obtiveram a Média mais baixa ( 41,2 ), e os da região Sul obtiveram a Média mais alta ( 48,4 ).
As demais médias foram : 47,2 na região Nordeste ; 44,7 na região Sudeste ; e 46,0 na região Centro-Oeste.
O Desvio padrão para o Brasil como um todo foi 12,7, sendo o maior Desvio padrão encontrado na região Norte ( 14,0 ) e o menor, na região Sudeste ( 11,7 ), indicando uma dispersão um pouco menor das notas desta última região.
A região que obteve a maior nota Máxima foi a Sul ( 92,6 ), ao passo que a região que atingiu a menor nota Máxima foi a Norte ( 70,1 ).
A Mediana do Brasil como um todo foi 46,0, sendo a maior Mediana obtida na região Sul ( 48,3 ), e a menor obtida no Norte ( 40,4 ).
A menor nota Mínima foi 5,0 na região Norte, enquanto a maior foi 20,1 na região Centro-Oeste.
Considerando -se as notas segundo Grande Região, observa -se que existe diferença estatisticamente significativa ao nível de 95 \% entre a menor Média, obtida na região Norte ( 41,2 ), e a maior Média, obtida na região Sul ( 48,4 ).
O comportamento das notas dos estudantes de todo o Brasil pode ser observado no Gráfico 3.1 que apresenta um histograma com a distribuição das mesmas.
Essa é uma distribuição unimodal com a moda no intervalo ( 40 ; 50 ].
Sudeste ( 11,7 ), indicando uma dispersão um pouco menor das notas desta última região.
A região que obteve a maior nota Máxima foi a Sul ( 92,6 ), ao passo que a região que atingiu a menor nota Máxima foi a Norte ( 70,1 ).
A Mediana do Brasil como um todo foi 46,0, sendo a maior Mediana obtida na região Sul ( 48,3 ), e a menor obtida no Norte ( 40,4 ).
A menor nota Mínima foi 5,0 na região Norte, enquanto a maior foi 20,1 na região Centro-Oeste.
Considerando -se as notas segundo Grande Região, observa -se que existe diferença estatisticamente significativa ao nível de 95 \% entre a menor Média, obtida na região Norte ( 41,2 ), e a maior Média, obtida na região Sul ( 48,4 ).
O comportamento das notas dos estudantes de todo o Brasil pode ser observado no Gráfico 3.1 que apresenta um histograma com a distribuição das mesmas.
Essa é uma distribuição unimodal com a moda no intervalo ( 40 ; 50 ].
A Tabela 3.3 apresenta informações referentes à Condição de Presença desagregadas por Categoria Administrativa e por Organização Acadêmica.
Da população total de Inscritos, 2.932 são de IES Privada e 136 de IES Públicas.
Em relação à Organização Acadêmica, a maior participação foi obtida por estudantes de Universidades ( 1.484 ), seguido de Centros Universitários ( 1.091 ) e Faculdades ( 493 ).
A Categoria Administrativa de maior abstenção foi a Pública ( 26,5 \% ), e entre as Organizações Acadêmicas foi a das Universidades ( 25,5 \% ), que ficou acima da abstenção nacional de 23,0 \%.
Já a Tabela 3.4 apresenta informações referentes à Média da nota final desagregadas por Categoria Administrativa e por Organização Acadêmica.
A Média das notas da prova como um todo foi 46,2.
Em relação à Categoria Administrativa, os estudantes das IES Públicas obtiveram Média mais alta ( 52,9 ), e os das IES Privadas obtiveram Média mais baixa ( 45,9 ), que a Média nacional.
Observa -se que existe diferença estatisticamente significativa entre as médias das notas das IES Públicas e Privadas.
A diferença entre as médias das regiões Sul e Norte ( 7,2 ), a maior e a menor Média, é superior à diferença entre IES Públicas e Privadas ( 7,0 ), caracterizando uma maior diversidade regional do que administrativa.
No tocante à Organização Acadêmica, apenas os Centros Universitários obtiveram Média mais alta ( 47,2 ) que a nacional.
As demais médias foram : 46,0 nas Universidades e 44,9 nas Faculdades.
O Desvio padrão para as IES Públicas ( 13,3 ) e para os Centros Universitários ( 13,1 ) foi superior ao do Brasil como um todo ( 12,7 ), indicando uma dispersão um pouco maior das notas nesta Categoria Administrativa e nesta Organização Acadêmica.
Constata -se que existe diferença estatisticamente significativa ao nível de 95 \% nas médias das notas dos estudantes dos Centros Universitários em relação aos das Faculdades. 3.1.2 Estatísticas Básicas no Componente de Formação Geral A Tabela 3.5 apresenta as Estatísticas Básicas em relação ao componente da prova que avalia a Formação Geral dos estudantes concluintes.
Os alunos de todo Brasil obtiveram desempenho médio de 50,4.
Quanto à variabilidade, o Desvio padrão das notas dos estudantes do Brasil, como um todo, foi 15,8.
A maior Média foi obtida na região Centro-Oeste ( 52,4 ), e a menor, na região Norte ( 46,1 ).
As demais médias foram : 51,2 na região Nordeste ; 49,7 na região Sudeste ; e 50,4 na região Sul.
Já o maior Desvio padrão foi obtido na região Norte ( 18,2 ), e o menor, na região Sudeste ( 14,9 ).
Os demais desvios padrões foram : 16,8 na região Nordeste, 15,4 da região Sul, e 17,0 na região Centro-Oeste.
A maior nota no Componente de Formação Geral da prova do Enade/2015 foi 95,4, obtida por, pelo menos, um aluno na região Centro-Oeste.
A menor nota Máxima foi obtida na região Sul ( 87,3 ).
A Mediana do Brasil, como um todo, foi 50,7, sendo a menor Mediana encontrada na região Norte ( 44,1 ) ; e a maior encontrada na região Centro-Oeste ( 52,7 ).
A nota Mínima nesta parte foi zero na região Sudeste.
Considerando -se as notas segundo Grande Região, observa -se que não existe diferença estatisticamente significativa entre as médias das notas do Componente de Formação Geral para cada uma das regiões do país. 3.1.2 Estatísticas Básicas no Componente de Formação Geral A Tabela 3.5 apresenta as Estatísticas Básicas em relação ao componente da prova que avalia a Formação Geral dos estudantes concluintes.
Os alunos de todo Brasil obtiveram desempenho médio de 50,4.
Quanto à variabilidade, o Desvio padrão das notas dos estudantes do Brasil, como um todo, foi 15,8.
A maior Média foi obtida na região Centro-Oeste ( 52,4 ), e a menor, na região Norte ( 46,1 ).
As demais médias foram : 51,2 na região Nordeste ; 49,7 na região Sudeste ; e 50,4 na região Sul.
Já o maior Desvio padrão foi obtido na região Norte ( 18,2 ), e o menor, na região Sudeste ( 14,9 ).
Os demais desvios padrões foram : 16,8 na região Nordeste, 15,4 da região Sul, e 17,0 na região Centro-Oeste.
A maior nota no Componente de Formação Geral da prova do Enade/2015 foi 95,4, obtida por, pelo menos, um aluno na região Centro-Oeste.
A menor nota Máxima foi obtida na região Sul ( 87,3 ).
A Mediana do Brasil, como um todo, foi 50,7, sendo a menor Mediana encontrada na região Norte ( 44,1 ) ; e a maior encontrada na região Centro-Oeste ( 52,7 ).
A nota Mínima nesta parte foi zero na região Sudeste.
Considerando -se as notas segundo Grande Região, observa -se que não existe diferença estatisticamente significativa entre as médias das notas do Componente de Formação Geral para cada uma das regiões do país.
O Gráfico 3.2 propicia a avaliação do desempenho dos estudantes do Componente de Formação Geral, a partir do histograma da distribuição das notas correspondentes.
A distribuição é unimodal, com moda no intervalo ( 40 ; 50 ], enquanto na prova, como um todo ( Gráfico 3.1 ), a moda foi alcançada apenas no intervalo ( 40 ; 50 ].
Nota -se que, no Gráfico 3.2, as notas apresentam maior dispersão do que no Gráfico 3.1 ( distribuição das notas da prova ), confirmado pela comparação dos desvios padrões : 12,7 para a nota da prova como um todo e 15,8 para o Componente de Formação Geral.
Na Tabela 3.6 são apresentadas as informações referentes ao desempenho dos concluintes do Componente de Formação Geral, em diferentes agregações : Categoria Administrativa e Organização Acadêmica.
Observa -se que existe diferença estatisticamente significativa entre as médias dos tipos de Categoria Administrativa.
A maior Média foi obtida por estudantes de IES Públicas ( 54,3 ), e a menor, por estudantes de IES Privadas ( 50,2 ).
Considerando -se o tipo de Organização Acadêmica, nota -se que não há diferença estatisticamente significativa entre as médias dos tipos de Organização Acadêmica. 3.1.3 Estatísticas Básicas do Componente de Conhecimento Específico A Tabela 3.7 apresenta as Estatísticas Básicas referentes ao Componente de Conhecimento Específico da área de Tecnologia em Design de Interiores.
A Média do desempenho dos alunos do Brasil, como um todo, foi 44,8.
A maior Média foi obtida na região Sul ( 47,7 ), e a menor, na região Norte ( 39,6 ).
As demais médias foram : 45,9 na região Nordeste ; 43,0 na região Sudeste ; e 43,9 na região Centro-Oeste.
Quanto à variabilidade das notas, o Desvio padrão do Brasil, como um todo, foi 13,9, sendo o maior Desvio padrão observado na região Norte ( 15,5 ), e o menor, na região Sudeste ( 13,0 ).
Os demais desvios foram : 14,5 na região Nordeste, 14,2 na região Sul e 14,0 na região Centro-Oeste.
A Mediana das notas dos estudantes de todo o Brasil foi 44,5.
A maior Mediana ocorreu na região Sul ( 47,0 ), e a menor, na região Norte ( 39,6 ).
As demais medianas foram : 45,7 na região Nordeste ; e 42,2 nas regiões Sudeste e Centro-Oeste.
A nota Máxima do Brasil, como um todo, foi 94,4, sendo obtida por, pelo menos, um aluno na região Sul.
As demais notas máximas foram : 66,9 na região Norte ; 83,4 na região Nordeste ; 76,8 na região Sudeste ; e 81,4 na região Centro-Oeste.
A menor nota Mínima foi obtida na região Nordeste ( 3,3 ) e a maior, na região Centro-Oeste ( 13,9 ).
Observa -se que existe diferença estatisticamente significativa da menor Média das notas do Componente de Conhecimento Específico da região Norte ( 39,6 ) em relação à maior Média, obtida na região Sul ( 47,7 ).
O intervalo de confiança da região Norte ( 8,7 ) é significativamente maior do que de todas as outras.
Assim como os Gráficos 3.1 e 3.2, o Gráfico 3.3, apresentado a seguir, permite uma avaliação do desempenho de concluintes de Tecnologia em Design de Interiores em relação ao Componente de Conhecimento Específico com um histograma da distribuição das notas correspondentes.
Dentre as três distribuições apresentadas, esta é a mais concentrada nas notas baixas.
Esta também é uma distribuição unimodal, e o grupo modal é o ( 40 ; 50 ], o mesmo grupo modal da prova como um todo e para a Formação Geral.
A Tabela 3.8 apresenta uma comparação dos resultados em relação à Categoria Administrativa e à Organização Acadêmica, agora levando em conta o desempenho de alunos do Componente de Conhecimento Específico da prova.
No que se refere à Organização Acadêmica, a maior Média foi dos Centros Universitários ( 45,9 ), vindo a seguir a das Universidades ( 44,6 ) e, depois, a das Faculdades ( 43,5 ).
O maior Desvio padrão, e acima do valor para o Brasil, como um todo, foi o dos Centros Universitários ( 14,4 ).
Os Centros Universitários também obtiveram as maiores notas Máxima ( 94,4 ) e Mediana ( 45,8 ).
Por sua vez, as Universidades obtiveram, respectivamente, 84,9 e 44,2 como nota Máxima e Mediana.
Já as Faculdades, obtiveram a nota Máxima ( 83,4 ) e a Mediana ( 43,2 ).
A menor nota Mínima foi 3,3 para as Universidades e a maior, foi 13,6 para as Faculdades.
Quanto à Categoria Administrativa, observa -se um comportamento semelhante àquele da prova como um todo e do Componente de Formação Geral, ou seja, existe 3.2 ANÁLISE DAS QUESTÕES OBJETIVAS Esta seção apresenta estatísticas selecionadas e histogramas das Questões Objetivas de Formação Geral ( 3.2.1 ).
São também apresentadas e comparadas as médias das subpopulações caracterizadas por Grande Região. 3.2.1 Componente de Formação Geral A Tabela 3.9 apresenta as Estatísticas Básicas relativas às oito questões objetivas do componente da prova que abrange a Formação Geral dos estudantes.
A Média do Brasil foi 47,9.
A menor Média foi encontrada na região Norte ( 42,9 ), e a maior, na região Centro-Oeste ( 50,9 ).
As demais médias foram : 49,7 na região Nordeste ; 46,6 na região Sudeste ; e 47,9 na região Sul.
O Desvio padrão do Brasil foi 19,4, sendo o maior Desvio padrão encontrado na região Norte ( 23,8 ), e o menor, na região Sul ( 18,3 ).
Os demais desvios foram : 20,3 na região Nordeste ; 18,9 na região Sudeste ; e 20,1 na região Centro-Oeste.
As medianas ( 50,0 ), as notas máximas ( 100,0 ) e as notas mínimas ( 0,0 ) foram iguais para todas as regiões, sem exceção. diferença estatisticamente significativa entre as médias das IES Públicas ( 52,4 ) e IES Privadas ( 44,5 ).
Observa -se que existem diferenças estatisticamente significativas ao nível de 95 \% no Componente de Conhecimento Específico entre as notas dos Centros Universitários em relação às das Faculdades. 3.2 ANÁLISE DAS QUESTÕES OBJETIVAS Esta seção apresenta estatísticas selecionadas e histogramas das Questões Objetivas de Formação Geral ( 3.2.1 ).
São também apresentadas e comparadas as médias das subpopulações caracterizadas por Grande Região. 3.2.1 Componente de Formação Geral A Tabela 3.9 apresenta as Estatísticas Básicas relativas às oito questões objetivas do componente da prova que abrange a Formação Geral dos estudantes.
A Média do Brasil foi 47,9.
A menor Média foi encontrada na região Norte ( 42,9 ), e a maior, na região Centro-Oeste ( 50,9 ).
As demais médias foram : 49,7 na região Nordeste ; 46,6 na região Sudeste ; e 47,9 na região Sul.
O Desvio padrão do Brasil foi 19,4, sendo o maior Desvio padrão encontrado na região Norte ( 23,8 ), e o menor, na região Sul ( 18,3 ).
Os demais desvios foram : 20,3 na região Nordeste ; 18,9 na região Sudeste ; e 20,1 na região Centro-Oeste.
As medianas ( 50,0 ), as notas máximas ( 100,0 ) e as notas mínimas ( 0,0 ) foram iguais para todas as regiões, sem exceção.
A Tabela 3.10 apresenta o Índice de Facilidade e o Índice de Discriminação ( ponto bisserial ) para cada uma das questões objetivas do Componente de Formação Geral.
Quanto ao Índice de Facilidade, foram usadas as seguintes cores para diferenciar o nível de dificuldade da questão :  Azul para as questões classificadas com índice Muito fácil ( \&gt ; =0,86 ), verde para as questões classificadas com índice Fácil ( 0,61 a 0,85 ), amarelo para as questões classificadas com Médio ( 0,41 a 0,60 ), vermelho para as questões classificadas com Difícil ( 0,16 a 0,40 ) e roxo para as questões classificadas com Muito difícil ( \&lt ; =0,15 ).
Já quanto ao Índice de Discriminação, foram usadas as seguintes cores para qualificar a questão :  As questões classificadas com índice Fraco receberam a cor vermelho ( \&lt ; =0,19 ), as classificadas com Médio receberam a cor amarelo ( 0,20 a 0,29 ), as classificadas com Bom receberam a cor verde ( 0,30 a 0,39 ) e as classificadas com Muito bom ( \&gt ; =0,40 ) receberam a cor azul.
As questões objetivas do Componente de Formação Geral, segundo o Índice de Facilidade, foram assim avaliadas : das oito questões, nenhuma teve o Índice de Facilidade classificado como Muito fácil.
Duas questões foram tidas como Fácil, por terem índice de acertos situado na faixa entre 0,61 e 0,85 ( de 61,0 \% a 85,0 \% de acertos ).
Três questões foram consideradas de dificuldade Médio, situando -se no intervalo entre 0,41 e 0,60 do Índice de Facilidade, ou seja, houve entre 41,0 \% e 60,0 \% de acertos, enquanto três questões foram classificada na categoria Difícil, situando -se no intervalo entre 0,16 e 0,40.
Por fim, nenhuma das questões apresentou menos de 15 \% de acertos, razão pela qual seria classificada como Muito difícil.
Como já comentado, para análise das questões objetivas, relativas à Formação Geral, segundo o poder de discriminação, utilizou -se o Índice de Discriminação ( ponto bisserial ). bisserial ).
Nesta análise, as questões foram assim avaliadas : cinco das oito questões apresentaram índice acima ou igual a 0,40 e, assim, foram classificadas com índice Muito bom para esse grupo de estudantes.
Três questões tiveram Índice de Discriminação Bom, com valor entre 0,30 e 0,39, para esse grupo de estudantes.
Nenhuma questão teve nível Médio ou Fraco de discriminação para esse grupo de estudantes.
O Índice de Facilidade variou de 0,27 a 0,71, e o de Discriminação, de 0,34 a 0,47.
As cinco questões com Índice de Discriminação Muito bom figuraram entre os diversos níveis de dificuldade desse conjunto : uma classificada na categoria Fácil ( questão 6 ) do Índice de Facilidade, duas, na categoria Médio ( questões 3 e 8 ) e duas, na categoria Difícil ( questões 2 e 7 ).
Em particular, a questão 2 foi a que apresentou o maior poder discriminatório, com índice 0,47, porém foi considerada difícil em termos de facilidade, com uma proporção de 0,35 acertos.
O máximo de acertos foi alcançado pela questão 5 com um Índice de Facilidade de 0,71.
A questão de número 3 apresentou um Índice de Facilidade de 0,52, ou seja, pouco mais da metade dos estudantes conseguiu resolvê -la, dentro do universo de participantes.
Seu Índice de Discriminação foi Muito bom ( 0,46 ).
Já a questão 6 obteve Índice de Discriminação Muito bom, 0,45, porém seu Índice de Facilidade foi Fácil ( 0,69 ).
O Gráfico 3.4, para exemplificar, analisa o comportamento da questão de número 2 de Formação Geral.
Trata -se de uma questão considerada mediana, com relação à facilidade e a que obteve o maior Índice de Discriminação dessa parte da prova.
Neste gráfico, cada uma das cinco curvas representa o percentual de respostas em determinada alternativa da questão, em função do número de acertos dos estudantes nessa parte da prova ( Formação Geral/Múltipla Escolha ), antes de possíveis eliminações pelo critério do ponto bisserial.
A curva em preto corresponde à alternativa C, a correta para esta questão.
Assim, observa -se que entre os estudantes com menor número de acertos, nessa parte do exame, a situação mais frequente foi a escolha da alternativa D ( em roxo ), incorreta.
Por exemplo, entre os estudantes que acertaram três questões, 35,0 \% escolheram a alternativa D, 20,4 \% a alternativa C ( correta ), 16,0 \% a alternativa A, 14,8 \% a alternativa B e 12,9 \% a alternativa E.
Entre os que acertaram três respostas entre as questões de múltipla escolha de formação geral, 0,8 \% deixou a questão em branco e 0,2 \% marcou mais de uma alternativa, invalidando a questão.
À medida em que o número de acertos aumenta, indicando desempenho melhor nessa parte da prova, aumenta concomitantemente a proporção de estudantes que selecionaram a alternativa correta C, atingindo 100 \% para os estudantes com 8 acertos.
Essa análise permite verificar como a questão discriminou os grupos de desempenho, justificando o alto índice obtido na questão.
Cumpre notar que não é possível inferir deste gráfico nem o índice de facilidade ( que seria uma média da proporção ponderada pela quantidade de alunos com cada uma das notas ), nem o índice de discriminação ponto bisserial, por razão equivalente.
No caso extremo no qual a grande concentração dos acertos dos alunos fosse abaixo de 4, o índice de facilidade seria obrigatoriamente abaixo de 40 \% ( neste exemplo ).
Caso a concentração fosse em 6 acertos ou mais, o índice seria obrigatoriamente acima de 60 \%.
Os gráficos relativos às demais questões de Formação Geral constam do Anexo I. 3.2.2 Componente de Conhecimento Específico A Tabela 3.11 apresenta as Estatísticas Básicas em relação às questões objetivas do Componente de Conhecimento Específico da prova de Tecnologia em Design de Interiores por Grande Região.
A Média do Brasil deste componente foi de 46,9.
A menor Média foi observada na região Norte ( 41,1 ), e a maior, na região Sul ( 49,4 ).
O Desvio padrão de todo o Brasil foi 14,8, sendo o menor Desvio padrão encontrado na região Sudeste ( 13,8 ), e o maior, na região Norte ( 17,8 ).
A Mediana de todo o Brasil foi 47,6, a mesma encontrada nas regiões Nordeste, Sul e Centro-Oeste.
Outras regiões apresentaram valores menores para a Mediana : 38,1 na região Norte e 42,9 na região Sudeste.
A nota Máxima da prova foi 95,2, obtida nas questões objetivas do Componente de Conhecimento Específico, por, pelo menos, um aluno na região Sul, e nas demais regiões, a nota Máxima da prova foi : 71,4 na região Norte ; 90,5 nas regiões Nordeste e Centro-Oeste ; e 81,0 na região Sudeste.
Nas regiões Norte e Nordeste a nota Mínima foi zero, enquanto a maior nota Mínima foi obtida na região Centro-Oeste ( 14,3 ).
A Tabela 3.12 apresenta os Índices de Facilidade e Discriminação ( ponto bisserial ) das questões objetivas do Componente de Conhecimento Específico para os estudantes de Tecnologia em Design de Interiores.
Para facilitar a diferenciação das questões, usaremos as mesmas cores da Tabela 3.10 para as diferentes classificações dos Índices de Facilidade e de Discriminação.
Dentre as questões objetivas da parte da prova relativa ao Componente de Conhecimento Específico, nenhuma questão foi anulada pela Comissão.
Desse modo, a classificação, quanto ao Índice de Facilidade, foi estabelecida com base em todas as 27 questões.
A partir dos índices obtidos, pode -se concluir que mais da metade das questões objetivas da prova foi considerada, pelo menos, Difícil : das 27 questões, 14 foram classificadas como Difícil ( a classificação modal ) ou como Muito difícil.
Quatro questões foram classificadas como Fácil, e outras nove consideradas como Médio.
Nenhuma questão foi considerada Muito fácil.
Já quanto aos índices de discriminação das questões objetivas do Componente de Conhecimento Específico da prova, tem -se como resultado a seguinte classificação : oito das 27 questões foram consideradas como boas, enquanto uma delas teve Índice de Discriminação Muito bom.
Assim, para a minoria das questões – nove em 27 – os Índices de Discriminação foram Bom ou Muito bom.
Dentre as demais, 12 delas foram classificadas como Médio e outras seis como Fraco, sendo 18, por conseguinte, a quantidade de questões nos dois patamares mais baixos de discriminação.
Constata -se, assim, que a prova – no que se refere ao Componente de Conhecimento Específico – possuía capacidade mediana de discriminar entre aqueles que dominam ou não o conteúdo.
Dentre as questões que alcançaram os maiores índices de discriminação, uma delas, a de número 34, foi classificada com Índice de Discriminação Muito bom, situando -se no intervalo de 0,40 ou mais do índice, sendo classificada na categoria A Tabela 3.12 apresenta os Índices de Facilidade e Discriminação ( ponto bisserial ) das questões objetivas do Componente de Conhecimento Específico para os estudantes de Tecnologia em Design de Interiores.
Para facilitar a diferenciação das questões, usaremos as mesmas cores da Tabela 3.10 para as diferentes classificações dos Índices de Facilidade e de Discriminação.
Dentre as questões objetivas da parte da prova relativa ao Componente de Conhecimento Específico, nenhuma questão foi anulada pela Comissão.
Desse modo, a classificação, quanto ao Índice de Facilidade, foi estabelecida com base em todas as 27 questões.
A partir dos índices obtidos, pode -se concluir que mais da metade das questões objetivas da prova foi considerada, pelo menos, Difícil : das 27 questões, 14 foram classificadas como Difícil ( a classificação modal ) ou como Muito difícil.
Quatro questões foram classificadas como Fácil, e outras nove consideradas como Médio.
Nenhuma questão foi considerada Muito fácil.
Já quanto aos índices de discriminação das questões objetivas do Componente de Conhecimento Específico da prova, tem -se como resultado a seguinte classificação : oito das 27 questões foram consideradas como boas, enquanto uma delas teve Índice de Discriminação Muito bom.
Assim, para a minoria das questões – nove em 27 – os Índices de Discriminação foram Bom ou Muito bom.
Dentre as demais, 12 delas foram classificadas como Médio e outras seis como Fraco, sendo 18, por conseguinte, a quantidade de questões nos dois patamares mais baixos de discriminação.
Constata -se, assim, que a prova – no que se refere ao Componente de Conhecimento Específico – possuía capacidade mediana de discriminar entre aqueles que dominam ou não o conteúdo.
Dentre as questões que alcançaram os maiores índices de discriminação, uma delas, a de número 34, foi classificada com Índice de Discriminação Muito bom, situando -se no intervalo de 0,40 ou mais do índice, sendo classificada na categoria Médio, quanto ao Índice de Facilidade.
A questão de número 13 foi a mais difícil dentre as 27 questões específicas, com baixo Índice de Facilidade, apenas 12,0 \% de acertos.
Essa questão apresentou poder discriminatório igualmente baixo, 0,12, o que comprova ter sido esta questão a mais difícil para os estudantes.
Destaca -se, também, a questão 27, com Índice de Facilidade 0,14, o que, em termos percentuais, corresponde a 14,0 \% de estudantes que responderam acertadamente.
Já 0,09 foi o seu Índice de Discriminação.
Tais questões foram, portanto, pelo critério ponto bisserial, consideradas inadequadas.
Por isso, as questões 13 e 27 foram eliminadas do cômputo da nota final.
Além destas duas, as demais questões com Índice de Discriminação Fraco, questões 20, 22, 25 e 31 também não foram utilizadas no cômputo final das notas, num total de seis questões eliminadas.
A título de exemplo das análises do comportamento das questões objetivas, o Gráfico 3.5 analisa a questão 34 do Componente de Conhecimento Específico.
Esta questão foi considerada pelos estudantes avaliados como uma questão de classificação Médio como nível de facilidade da prova, apresentou Índice de Facilidade 0,54, ou seja, 54,0 \% dos estudantes assinalaram acertadamente a opção C, correspondente ao gabarito.
Seu Índice de Discriminação foi igual a 0,40, classificado como Muito bom, também sendo esta questão a que apresentou o maior índice discriminatório.
Neste gráfico, cada uma das cinco curvas representa o percentual de respostas em determinada alternativa da questão 34, em função do número de acertos dos estudantes nessa parte da prova, antes de possíveis eliminações de questões pelo critério do ponto bisserial.
A alternativa correta C, representada no gráfico pela curva em preto, foi escolhida em maiores proporções pelos alunos com desempenho melhor nessa parte da prova.
Já as alternativas incorretas, também denominadas distratores, foram selecionadas, principalmente, por aqueles com notas mais baixas.
Neste caso também a soma não é sempre 100 \% por causa das questões não respondidas ou com mais de uma opção marcada.
Aqueles com nota zero, na sua quase totalidade deixaram esta questão em branco ou marcaram mais de uma alternativa, comportamento considerado inválido.
A proporção de alunos que selecionou a resposta correta C aumenta gradativamente, chegando a atingir 100 \% para 19 acertos ou mais, enquanto a proporção dos que escolheram alternativas incorretas decai, a partir do primeiro acerto, como função do número de acertos nessa parte da prova.
Os gráficos relativos às demais questões do Conhecimento Específico constam do Anexo I. 3.3 ANÁLISE DAS QUESTÕES DISCURSIVAS Esta seção apresenta estatísticas selecionadas e histogramas das Questões Discursivas de Formação Geral ( 3.3.1 ) e Conhecimento Específico ( 3.3.2 ).
São também apresentadas e comparadas as médias de subpopulações caracterizadas por Grande Região. 3.3.1 Componente de Formação Geral As análises dos resultados de desempenho dos estudantes de Tecnologia em Design de Interiores nas duas questões discursivas relativas à Formação Geral encontram -se na Tabela 3.13 e no Gráfico 3.6.
Na Tabela 3.13, observa -se que a nota Média nesse conjunto de questões foi próxima da obtida nas objetivas.
Os estudantes, de todo o Brasil, obtiveram, em Formação Geral, Média 47,9 nas questões objetivas e 54,1 nas questões discursivas.
No entanto, pode -se notar um aumento do Desvio padrão de 19,4 nas questões objetivas do Componente de Formação Geral dos alunos de todo o Brasil, para 22,4 nas questões discursivas do mesmo componente.
A maior Média foi obtida na região Centro-Oeste ( 54,6 ), e a menor, na região Norte ( 51,0 ).
A Mediana de todo o Brasil, neste componente, foi 60,5, nas regiões Nordeste e Sul a Mediana foi maior, 61,5 em ambas.
Nas demais regiões foi menor : Norte ( 57,5 ), Sudeste ( 59,5 ) e Centro-Oeste ( 60,0 ).
A nota Máxima ( 95,0 ) foi obtida na região Nordeste, enquanto a menor nota Máxima foi obtida na região Norte ( 87,5 ).
A nota Mínima ( 0,0 ) foi a mesma em todas as regiões do Brasil, sem exceção.
O Gráfico 3.6 representa a distribuição das notas nas questões discursivas do Componente de Formação Geral.
A moda desta distribuição ocorre no intervalo ( 60 ; 70 ].
Destaca -se, também, o intervalo [ 0 ; 10 ], um máximo local com distribuição próxima a 10 \% do total de notas, sendo que no intervalo [ 0 ; 10 ] se inclui, além da nota zero, a frequência de alunos que deixaram este tipo de questão em branco.
Na sequência, os resultados verificados para cada uma das questões discursivas de Formação Geral serão apresentados, estabelecendo -se relações com os conteúdos abordados em cada uma delas.
Os comentários da Banca de docentes corretores, a respeito do observado na correção das respostas dos estudantes, suas impressões e conclusões serão apresentados junto à análise de cada questão.
Cumpre esclarecer que, tendo em vista que as questões discursivas de Formação Geral são padronizadas, ou seja, constam de todas as provas, os comentários da Banca são os mesmos para todas as carreiras acadêmicas, sendo direcionados a todos os estudantes que participaram do Enade/2015.
A seguir, serão analisados os desempenhos da Área de Tecnologia em Design de Interiores nas duas questões discursivas de Formação Geral do Enade/2015, comparando -se os resultados obtidos com comentários para cada questão. 3.3.1.1 Análise de Conteúdo da Questão Discursiva 1 do Componente de Formação Geral Os dados de Tecnologia em Design de Interiores, obtidos a partir das respostas à questão 1, encontram -se na Tabela 3.14 e no Gráfico 3.7.
Nessa questão – de melhor desempenho dentre as duas de Formação Geral – os alunos, de todo o Brasil, tiveram Média 59,6.
A maior Média para a questão 1 foi obtida na região Centro-Oeste ( 61,2 ), e a menor, na região Norte ( 56,5 ).
Quanto à variabilidade das notas, o Desvio padrão de todo o Brasil foi 24,9.
O menor Desvio padrão foi obtido na região Sudeste ( 23,4 ), e o maior Desvio padrão foi obtido na região Norte ( 28,2 ).
As medianas ( 65,0 ), as notas máximas ( 95,0 ) e as notas mínimas ( 0,0 ) foram iguais para todas as regiões, exceto a nota Máxima na região Norte ( 90,0 ).
O Gráfico 3.7 mostra a distribuição das notas na questão discursiva 1 do Componente de Formação Geral.
Observa -se que a maior frequência corresponde aos alunos que obtiveram nota dentro do intervalo ( 70 ; 80 ], correspondendo à moda da distribuição.
Destacam -se, também, os alunos que deixaram a questão em branco, representando um máximo local com mais de 10 \% do total. 3.3.1.2 Comentários sobre a correção de Conteúdo das respostas à Questão Discursiva 1 O enunciado da Discursiva 1 de Formação Geral estava claro e indicava um tema de extrema relevância, qual seja : a busca pela educação universal, especificamente, que apontava um caso onde ocorreu a discriminação de acesso ao ensino formal atingindo uma menina, representando a categoria das mulheres discriminadas em tantos campos da vida e, mais ainda, em sociedades que cultuam o poder soberano dos homens.
A linguagem utilizada foi clara, indicando a existência dos fatos vinculados à luta da menina Malala pelo acesso ao ensino formal das jovens do seu país, e o confronto ideológico que gerou o ataque à sua integridade física.
Diante da repercussão internacional, houve o reconhecimento pelo prêmio Nobel, e a ONU resolveu apoiá -la como símbolo da luta pela igualdade entre os sexos buscando a universalização da educação.
O conteúdo está integrado às Diretrizes, vez que o acesso ao ensino formal está vinculado ao direito constitucional à educação, que, também, é protegido pelo Estatuto da Criança e do Adolescente e nos pactos de Direitos Humanos aderidos pelo Brasil.
Por outro lado, o texto demandava uma reflexão sobre a realidade dos indivíduos que não têm acesso à educação e a situação das mulheres excluídas desse processo por segmentos da sociedade.
Como se tratava de texto contendo uma ampla gama de possibilidades, era natural a ocorrência de interpretações diversas das esperadas.
Releva notar que o enunciado pretendeu, primordialmente, identificar a questão de gênero, apontando para as dificuldades que as mulheres possuem, não somente no Brasil, mas em vários países do mundo.
Ao analisar as provas realizadas, constatamos que no item ‘ a ’ poucos foram os que identificaram os estatutos normativos indicados no padrão de resposta.
Essa ausência pode ser justificada tendo em vista que, no enunciado indicado, não existia elemento que pudesse ser utilizado pelo concluinte para embasar sua resposta nesse caminho.
As respostas foram mais gerais, identificando que a educação constitui um item fundamental para o desenvolvimento do indivíduo e que as mulheres sofrem mais com as limitações impostas por países com regime patriarcal, que discriminam o sexo feminino e, com frequência, proíbem as mulheres de frequentar a escola.
Deve -se ter em vista que a indicação de resposta que vinculava o texto aos estatutos normativos representou um raciocínio próprio dos estudantes de Direito que, naturalmente, encaminharam suas respostas para a fundamentação nessa linha identificando, não somente a Constituição, mas também leis, regulamentos, etc.
Tal linha de argumentação não foi objeto da formação dos estudantes das demais áreas.
A maior parte dos concluintes estabeleceu o tema educação formal associado à busca pelo conhecimento, como acesso ao ensino, indicando a independência intelectual e o desenvolvimento da consciência crítica, limitando as manipulações e permitindo amplo acesso ao mercado de trabalho, traduzindo uma libertação econômica e ascensão social.
No tocante ao item ‘ b ’, restou claro que a maioria esmagadora das respostas confluiu para as indicações do padrão apresentado, sendo relevante o tema de igualdade de gênero, além de candente para as concluintes do sexo feminino que desenvolveram os textos de forma concatenada, fundamentada e, muitas vezes, lamentando que, nos dias atuais, ainda se discuta a necessidade de igualar os gêneros na sociedade.
Nessa linha, foram identificadas respostas, consoante com o padrão apresentado, especialmente nos seguintes temas :  violência física e psicológica contra a mulher, incluindo a Lei Maria da Penha, no caso específico do Brasil ;  tolerância/intolerância a vestimentas, trajes, comportamentos socialmente estereotipados ;  aspectos socioculturais que impõem à mulher uma condição de submissão na sociedade, tais como a impossibilidade de a mulher manifestar seus desejos e posicionamentos em algumas culturas, entre outros ;  igualdade/desigualdade de gênero, por exemplo no mercado de trabalho, em relação à desigualdade salarial ;  ideais de Liberdade, Igualdade e \textbf{Fraternidade} : referência a esses ideais como possibilidade de equilibrar as relações de poder entre homens e mulheres.
Em muitos casos, também, os concluintes fizeram um liame entre o acesso à educação e a melhora das relações entre homens e mulheres ao permitir que a conclusão de maior grau de educação formal permita o acesso a cargos públicos relevantes, empregos bem remunerados e posições de poder, tanto na área pública quanto na área privada.
Como as opções do segundo item foram mais amplas, o padrão de resposta identificou as linhas de pensamento adotadas pelos concluintes.
A maior parte das respostas identificou, claramente, o problema da relação de gênero e suas consequências na sociedade.
Também foi frequente constatarem que existe efetiva melhora na situação das mulheres no Brasil, mas que tal movimento não ocorre em todo o mundo.
Aliás, o exemplo de Malala conduziu a uma reflexão em âmbito mundial, e os concluintes apresentaram suas análises críticas à situação da mulher em países que não admitem sequer a sua inserção no plano educacional e, muito menos, no mercado de trabalho.
A maior parte dos concluintes apresentou concatenação lógica nas suas respostas e uma minoria, pouco expressiva, não demonstrou esse domínio de desenvolvimento lógico dos textos.
Ressalte -se que essa referência não inclui a utilização do padrão culto da língua portuguesa que tem avaliação própria. desenvolvimento lógico dos textos.
Ressalte -se que essa referência não inclui a utilização do padrão culto da língua portuguesa que tem avaliação própria. 3.3.1.3 Análise de Conteúdo da Questão Discursiva 2 do Componente de Formação Geral A Tabela 3.15 mostra que o desempenho médio dos estudantes na questão discursiva 2 ( média 52,3 ) foi inferior ao obtido na questão discursiva 1 ( média 59,6 ).
As regiões Sudeste e Sul foram aquelas cuja Média, nessa questão, foi maior ( 52,7 ), e a de menor Média foi a região Norte ( 48,6 ).
Quanto à variabilidade das notas, o Desvio padrão de todo o Brasil foi 25,6, superior ao obtido na questão discursiva 1 ( 24,9 ).
O maior desvio nessa questão foi obtido na região Nordeste ( 27,8 ), enquanto o menor foi obtido na região Sudeste ( 24,1 ).
A Mediana de todo o Brasil foi 60,0, a mesma das regiões Nordeste, Sudeste e Sul, enquanto nas regiões Norte ( 56,0 ) e Centro-Oeste ( 59,0 ) a Mediana foi menor.
A nota Máxima ( 100,0 ) foi obtida na região Nordeste.
A nota Mínima ( 0,0 ) foi a mesma em todas as regiões do Brasil, sem exceção.
O Gráfico 3.8 mostra a distribuição das notas na questão discursiva 2 do Componente de Formação Geral.
Observa -se que a maior frequência corresponde aos alunos que obtiveram nota situada no intervalo ( 60 ; 70 ].
Nota -se, ainda, que as notas ficaram mais distribuídas ao longo dos intervalos em comparação à questão discursiva de número 1. 3.3.1.4 Comentários sobre a correção de Conteúdo das respostas à Questão Discursiva 2 O enunciado da Discursiva 2 estava, igualmente, claro, indicando um tema de extrema relevância, qual seja : o reconhecimento de manifestações artísticas nascidas na periferia, como integrante da cultura reconhecida pela sociedade.
A linguagem utilizada foi clara, indicando a análise dos fatos e suas consequências no cotidiano dos partícipes e o confronto com outras visões culturais que são consideradas mais adequadas pela sociedade civil.
O conteúdo está integrado às Diretrizes, vez que o aspecto cultural e sua inter-relação com o conhecimento haurido no nível superior é essencial para a compreensão de mundo do universitário concluinte, com reflexo na sua atuação profissional.
O texto demanda uma reflexão sobre a realidade dos indivíduos que não têm acesso aos itens básicos da sociedade, como educação, saúde, cultura tradicional e segurança e que, apesar disso, criam alternativas para o seu lazer.
A incorporação do funk como manifestação cultural e artística da sociedade brasileira encontrou, basicamente, duas correntes de resposta : • os que aquiescem positivamente à integração ; • os que rejeitam, acerbamente, tal possibilidade.
No meio, uma corrente que somente admite a integração com modificações na conduta dos integrantes do movimento, com a exclusão de letras exclusivamente com apelos sexuais e incitação à violência.
A interpretação positiva indica a origem da música na periferia, notadamente nas comunidades cariocas e sua projeção para o Brasil e o mundo.
Nessa perspectiva, surge o reconhecimento como cultura popular comparada ao samba, axé e outros ritmos também nascidos no Brasil.
Também ocorre comparação com o movimento hip-hop americano.
Já a análise negativa prioriza a vinculação do ritmo com traficantes, prostitutas, vilipêndio às mulheres, falta de controle quanto ao som, local dos eventos, transações sexuais em público, uso frequente de drogas, violência \textbf{exacerbada}.
Faltam, para essa corrente, elementos básicos da cultura, tratando -se de música sem autonomia como movimento de mudança ou aspecto próprio.
A corrente intermediária reconhece o vínculo originário com as comunidades dominadas por quadrilhas de traficantes, mas admite o seu reconhecimento de acordo com a mudança dos atores que transitam na modalidade, como músicos, \textbf{intérpretes}, para adequar as letras das canções expurgando os aspectos considerados \textbf{nocivos} como o apelo sexual, as indicações positivas para a violência e apologia às drogas.
Um aspecto bastante criticado nas músicas dessa modalidade foi a exploração da mulher, entendida como um objeto de satisfação masculina.
É relevante tal constatação, pois a questão pertinente à Malala reflete a luta pela valorização da mulher.
Os concluintes compreenderam a ideia que foi problematizada e escolheram uma estratégia adequada para solucioná -la.
O caráter dinâmico da cultura foi comumente expressado, não através do enfrentamento abordado no texto, mas mediante a indicação ( de forma rasa e genérica ) das dificuldades e consequente aceitação social de manifestações culturais, antes marginalizadas e que hoje integram o acervo cultural nacional, tais como : samba, grafite, capoeira e forró.
Quanto às reflexões de Laraia, a grande maioria reconheceu o funk como manifestação cultural, porém, o identificou mais na sua expressão musical, no balanço de seu ritmo, tecendo contundentes críticas, e mesmo forte discordância quanto ao conteúdo e mensagem das letras quando remetem à depreciação da imagem feminina ou induzem à apologia ao crime.
De resto, os textos abordaram bem mais o item dois do padrão de Respostas.
Uma minoria enfrentou os temas mudança e dinamismo cultural.
As respostas da corrente positiva indicam que o preconceito sempre acompanhou as manifestações populares como o samba, na sua origem.
Admitem que diante da miscigenação surgem movimentos culturais de várias naturezas, o que inclui o funk.
Também indicam a dificuldade de lazer dos membros das comunidades o que faz surgirem eventos para preencher essa necessidade.
Nesse âmbito, apontam a música como alternativa de ascensão social e integração aos bens de consumo.
As respostas da corrente negativa vinculam a música ao tráfico de drogas cujos elementos comandam as comunidades, propiciando que os bailes sejam mecanismos de venda e consumo dessa mercadoria ilícita.
Na mesma toada, indicam a vinculação com a exploração sexual, principalmente das mulheres maiores e menores de idade.
Indicam aspectos imorais nas letras.
Apresentam os locais dos bailes como inadequados e provocadores da \textbf{desordem} pública.
Recusam que músicas com discriminação às mulheres sejam consideradas como cultura.
A corrente intermediária indica a possibilidade de adoção do funk como cultura desde que ocorram modificações nas letras e no comportamento dos atores ( músicos, letristas e público ), retirando os aspectos negativos já apontados.
Fala -se em “ funk do bem ”.
A maior parte das respostas apresentadas indica que os concluintes têm conhecimento dos conflitos sociais e, no caso da questão do funk, conseguem identificar a origem popular da música.
A partir daí é que ocorre a cisão entre os que defendem o movimento como cultura, aduzindo a ascensão social dos artistas envolvidos dos que vinculam a música com atos ilícitos, rejeitando sua inserção no âmbito cultural.
Em algumas provas, há dificuldade de concatenação lógica de raciocínio, mas a maior parte consegue desenvolver o texto de forma razoável.
Essa questão permite visualizar a divisão social com os pensamentos \textbf{liberais} e conservadores, sendo bastante claros nas respostas dos concluintes. 3.3.1.5 Análise de Língua Portuguesa das Questões Discursivas do Componente de Formação Geral Os dados de Tecnologia em Design de Interiores, obtidos a partir das respostas às questões discursivas do Componente de Formação Geral, no que tange à Língua Portuguesa, encontram -se na Tabela 3.16 e no Gráfico 3.9.
Nesse aspecto, os alunos, de todo o Brasil, tiveram Média 53,5.
A maior Média com respeito à Língua Portuguesa foi obtida na região Centro-Oeste ( 54,1 ), e a menor, na região Norte ( 50,2 ).
Quanto à variabilidade das notas, o Desvio padrão de todo o Brasil foi 22,8.
O menor Desvio padrão foi obtido na região Sudeste ( 21,3 ) e o maior Desvio padrão foi obtido na região Norte ( 25,1 ).
A Mediana das notas de Língua Portuguesa foi 59,5 para o Brasil como um todo, enquanto a maior Mediana e a menor Mediana foram obtidas, respectivamente nas regiões Nordeste ( 60,5 ) e Norte ( 57,5 ).
A nota Máxima para todo o Brasil foi de 97,5, com, pelo menos, um aluno tirando essa nota nas regiões Nordeste e Sul.
Além disso, a nota Mínima foi zero em todas as regiões do país, sem exceção.
O Gráfico 3.9 mostra a distribuição das notas de Língua Portuguesa do Componente de Formação Geral.
Observa -se que a maior frequência ( cerca de 25 \% ) corresponde aos alunos que obtiveram nota no intervalo ( 60 ; 70 ]. 3.3.1.6 Comentários sobre a correção das respostas de Formação Geral com respeito à Língua Portuguesa Ao encaminhar as questões 1 e 2 na direção da produção de um texto dissertativo, espera -se que o participante utilize seus conhecimentos sobre o assunto e estruture seus textos de acordo com as exigências do registro formal próprio dessa situação comunicativa.
Essa configuração determina exigências quanto : à adequação da seleção vocabular, ao desenvolvimento do conteúdo, à estruturação sintática dos períodos, à organização lógica das ideias, à utilização de procedimentos de encadeamento textual e referenciação, à obediência às exigências morfossintáticas próprias da modalidade escrita da norma-padrão, ao respeito às regras ortográficas e às regras de acentuação gráfica.
O padrão de resposta utilizado na avaliação das questões 1 e 2 considerou os aspectos relevantes ao bom desempenho linguístico como competências distintas, de modo a permitir um mapeamento detalhado do domínio dos recursos disponíveis na Língua Portuguesa para a comunicação escrita formal.
Com base nesse objetivo, foram avaliados os seguintes aspectos : a ) Estruturação textual condizente com o gênero solicitado ( texto dissertativo ) e o modo de organização textual expositivo adequado ao gênero – essa competência envolve :  distribuição do conteúdo do texto em parágrafos, de modo a garantir a sua organização temática ;  estruturação sintática condizente com o padrão da modalidade escrita formal da língua portuguesa de modo a garantir a clareza necessária ;  utilização de operadores discursivos que contribuam para a progressão temática do texto, estabelecendo relações lógicas entre as ideias apresentadas, tanto do ponto de vista intrafrasal, como do interfrasal ;  utilização de procedimentos de referenciação lexical e pronominal que permitam a retomada de referentes textuais ;  utilização de sinais de pontuação que contribuam para a organização lógica da frase e do texto.
Espera -se, portanto, que o participante recorra a procedimentos linguístico-discursivos para organizar seu texto, permitindo o encadeamento lógico entre suas partes de forma a garantir a progressão e a coerência textual.
Isso significa que os seguintes procedimentos foram penalizados, de acordo com o padrão de resposta proposto :  elaboração de frases fragmentadas que comprometam a estrutura lógico-gramatical do texto ;  sequência justaposta de ideias sem encaixamentos sintáticos, reproduzindo hábitos da oralidade ;  elaboração de frase com apenas oração subordinada, sem oração principal ;  emprego equivocado do conector ( preposição, conjunção, pronome relativo, alguns advérbios e locuções adverbiais ) que não expresse a relação lógica adequada entre dois trechos do texto e prejudique a compreensão da mensagem ;  repetição ou substituição inadequada de palavras sem se valer dos recursos oferecidos pela língua ( pronome, advérbio, artigo, sinônimo ) ;  utilização inadequada dos sinais de pontuação comprometendo a clareza textual. b ) Respeito às convenções ortográficas da norma-padrão da Língua Portuguesa – essa competência envolve o domínio das regras de acentuação gráfica e da grafia padrão das palavras ( com ausência de abreviaturas próprias da linguagem da internet ), de acordo com as convenções estabelecidas pela legislação em vigor e consubstanciadas no Vocabulário Ortográfico da Língua Portuguesa, editado pela Academia Brasileira de Letras ( com aceitação da legislação anterior, no caso das regras relativas ao uso do hífen e da acentuação gráfica, já que o exame foi realizado no mês de \textbf{novembro} de 2015 ).
A competência relativa ao respeito às convenções ortográficas envolve : • grafia correta das palavras ; • respeito às regras de acentuação gráfica ; • emprego de maiúsculas em início de frase, em nomes próprios de pessoas, lugares ou instituições ; • ausência de abreviações como p/, vc, tb, pq, tá, né, usadas muitas vezes em escrita informal e na internet ; • obediência às regras de separação de sílabas no final da linha. c ) Domínio dos diferentes aspectos morfossintáticos próprios da modalidade escrita formal da norma-padrão da Língua Portuguesa – os itens relacionados a essa competência são : a concordância nominal, a concordância verbal, a regência nominal, a regência verbal, a flexão nominal, a flexão verbal, a correlação entre os tempos verbais, a colocação pronominal.
O domínio dessa competência envolve : • flexão do verbo para estabelecer concordância de número com o sujeito da frase ; • flexão do artigo, do adjetivo e do pronome para concordar em número e em gênero com o substantivo a que se referem ; • respeito à regência nominal e à verbal, utilizando a preposição adequada depois de um substantivo, de um verbo ou de um adjetivo ; emprego da preposição antecedendo o pronome relativo e emprego adequadamente os pronomes relativos “ cujo(a) ” e “ onde ” ; • emprego adequado do acento grave indicador da crase entre uma preposição e um artigo ( a+a ) ; • atendimento às regras de colocação pronominal ( próclise e ênclise ), distintas dos hábitos da oralidade ou da escrita informal. d ) Seleção vocabular adequada à modalidade escrita formal da Língua Portuguesa, exigida pela situação comunicativa – essa competência envolve a precisão na utilização do vocabulário relacionado à temática solicitada pela questão ; a ausência de marcas da oralidade, como termos de sentido muito genérico ( “ coisa ”, “ negócio ”, “ você ” ) e termos de registros mais informais ( como gírias, jargões, frases feitas, ditados populares, termos regionais ).
Assim, espera -se que o participante respeite a adequação vocabular não empregando gírias ou expressões coloquiais, evite repetição desnecessária de palavras e utilize um vocabulário mais formal, como solicitado por um texto dissertativo.
A escolha dessas competências para subsidiar o processo de avaliação apoia -se na concepção de que, no desempenho dos graduandos, a modalidade escrita tem apresentado uma intensa simplificação, originada no padrão da modalidade oral da Língua Portuguesa.
No caso do texto de base dissertativa, inscrito em um registro formal, a distância entre as duas modalidades é ainda maior, o que provoca situações de hipercorreção ( desvios provocados pela utilização inadequada de uma regra da norma-padrão ) e de truncamentos sintáticos ( estruturas frasais incompreensíveis devido à complexidade sintática própria da modalidade escrita ).
Observam -se, então, os seguintes aspectos que marcam essa distinção entre as duas modalidades, devido à excessiva simplificação da modalidade falada : a ) redução drástica de estruturas subordinadas, compensada pelo aumento na frequência de estruturas coordenadas e absolutas, por um lado, ou pela elaboração de estruturas truncadas pelo excesso de ideias sem a devida conexão subordinativa ; b ) redução no uso de operadores argumentativas para expressar relações lógicas essenciais à construção do texto, substituídas pela exigência de inferência por parte do interlocutor para suprir a sua ausência ; c ) redução do uso do subjuntivo, ao lado da ampliação do uso do indicativo combinado a estruturas frasais coordenadas ou absolutas ; d ) empobrecimento do processo de referenciação, com a repetição de pronomes ou nomes ; e ) simplificação extrema da marcação da categoria tempo na morfologia verbal ; f ) falta de domínio de vocabulário mais abstrato e de maior complexidade, essencial ao desenvolvimento do processo dissertativo ; g ) redução drástica no emprego da acentuação gráfica, processo intensificado pela intensa utilização das redes sociais.
Os aspectos macroestruturais da elaboração do texto não foram avaliados neste processo porque dizem respeito à avaliação do conteúdo.
Portanto, a banca de Formação Geral, composta por profissionais de diferentes áreas do conhecimento, se encarregou da avaliação do atendimento ao solicitado no enunciado das questões do ponto de vista do desenvolvimento do conteúdo.
A avaliação do desempenho linguístico considerou, portanto, três grandes grupos de competências, segundo os aspectos explicitados anteriormente : a ) domínio das convenções ortográficas : grafia de vogais e consoantes, uso de maiúsculas e minúsculas, emprego do hífen e acentuação gráfica ; b ) domínio dos procedimentos de estruturação textual do ponto de vista microestrutural : organização interna dos períodos, emprego de conectores para a articulação lógica entre os períodos e entre os parágrafos, emprego de marcas de referenciação lexical e pronominal ; utilização dos sinais de pontuação que contribuem para a organização lógica da frase. c ) domínio das regras de caráter morfossintático estabelecidas como modelares do ponto de vista da modalidade escrita formal da norma-padrão da Língua Portuguesa : concordância nominal e verbal, regência nominal e verbal, colocação pronominal, flexão nominal e verbal, correlação entre tempos e modos verbais, ausência de marcas de oralidade.
A seleção vocabular adequada à modalidade escrita formal da Língua Portuguesa foi incorporada à essa última competência, tendo em vista a intersecção entre as duas do ponto de vista das exigências do registro formal da modalidade escrita da norma-padrão.
Apresenta -se, a seguir, o padrão de resposta aprovado pelo Inep e já utilizado na avaliação do desempenho linguístico das questões 1 e 2 de Formação Geral no Enade/2013 e no Enade/2014.
Às competências, reunidas nos três grupos descritos anteriormente, foram atribuídos os seguintes \textbf{pesos} relativos : aspectos ortográficos ( 20 \% ) ; aspectos textuais ( 40 \% ), aspectos morfossintáticos e vocabulares ( 40 \% ).
Com base na avaliação do processo desenvolvido nos dois anos anteriores, utilizaram -se quatro níveis para a avaliação do desempenho linguístico nas três competências.
O nível zero ficou reservado para casos especiais, como fuga total ao tema, palavras soltas, produção de uma frase incompleta, entre outros.
A avaliação das respostas elaboradas pelos participantes revelou um resultado coerente na comparação entre as duas questões.
Entretanto, em relação aos resultados do Enade/2014, observou -se uma melhora do domínio de todas as competências, talvez devido ao perfil dos cursos analisados, como Direito e Jornalismo, que pressupõem a utilização profissional da linguagem.
Observou -se uma diferença no desempenho dos participantes em relação às duas questões discursivas, provavelmente em função dos temas solicitados : - a Questão 1 solicitava que o participante desenvolvesse um texto dissertativo sobre o significado da premiação de Malala, abordando dois aspectos : o direito das jovens à educação formal e as relações de poder entre homens e mulheres no mundo.
Esse encaminhamento favoreceu o desenvolvimento mais articulado do texto, o que se refletiu em melhor desempenho dos aspectos vocabulares e textuais. - a Questão 2 solicitava que o participante desenvolvesse um texto dissertativo para discutir a questão do reconhecimento do funk como legítima manifestação artística e cultural do povo brasileiro.
Por ser um tema muito polêmico, o tema do preconceito sócio-cultural-econômico favoreceu a construção de textos mais críticos, dividindo o conjunto de participantes em dois grupos : os que defendiam esse tipo de manifestação artística como legítima e os que a criticavam de maneira veemente, chegando, inclusive, a utilizar muitos termos de baixo calão.
Essa última atitude teve como consequência a elaboração de textos mal desenvolvidos do ponto de vista sintático-discursivo, com muitos truncamentos.
Os resultados revelaramm essa diferença.
Ao analisar os dados de cada questão, constata -se que a Questão 1 teve maior número de questões de médio ou bom desempenho, em relação à Questão 2.
Por outro lado, a Questão 2 apresentou um índice maior de notas baixas em relação à Questão 1.
Os resultados da avaliação correspondem aos seguintes aspectos observados em cada competência : a ) aspectos ortográficos : o desempenho dos participantes revelou uma diferença muito grande nos dois aspectos analisados : baixo índice de desvios ortográficos e grande índice de desvios de acentuação.
Em vários casos, ocorre ausência completa de acentuação gráfica.
Os resultados revelam que a tendência dominante entre os universitários brasileiros é a eliminação da acentuação gráfica, talvez motivada pelos hábitos relacionados às redes sociais e pela ausência de esclarecimento dos meios de comunicação, das autoridades e das escolas sobre as decisões do Acordo Ortográfico da Língua Portuguesa de 1990.
Os casos mais sistemáticos de eliminação do acento indicador da sílaba tônica são :  palavras proparoxítonas ( “ proximos ”, “ politicos ” ) ;  palavras paroxítonas terminadas em ditongo crescente ( “ necessario ”, “ noticiarios ”, “ individuo ”, “ dependencia ”, “ varios ”, “ propria ” ) ;  palavras oxítonas ( “ ninguem ”, “ esta ”, “ ate ”, “ tambem ”, “ alvara ” ) ;  uso indevido do acento gráfico em palavras como “ genêro ”, “ melâncolia ”, por exemplo.
Quanto ao domínio das convenções relativas à grafia das palavras, observam -se desvios como : I. hipercorreção pela escolha de “ e ” no lugar de “ i ” ou de “ i ” no lugar de “ e ”, por influência de hábitos da oralidade ( “ descriminação ” por “ discriminação ”, “ entervenção ” por “ intervenção ”, “ esteriotipada ” por “ estereotipada ” ) ; II. eliminação do “ r ” marcador do infinitivo verbal ( “ esta ” no lugar de “ estar ”, “ estuda ” no lugar de “ estudar ” ).
III. eliminação de uma sílaba na palavra “ educação ”, que foi grafada “ edução ” por um grande número de participantes.
IV. outros casos de desvios de grafia relacionados à variação diastrática podem ser observados em “ estrupo ”, “ subjulgadas ”, “ precoseito ”, “ soubre ”, “ vecendo ”, “ indesencia ”, “ apolojia ”, “ fulga ”.
V. desvios de segmentação podem ser atribuídos à confusão entre o \textbf{morfema} gramatical “ mos ” e o pronome oblíquo “ nos ”, em “ chegar mos ”, “ enfrentar -mos ”.
VI. inadequação no uso da maiúscula para destacar determinadas palavras-chave do texto, como “ Violência ”, “ Brasileiros ”.
Destaque -se, também, o grande número de participantes que grafam os textos inteiramente em caixa alta.
Vale observar, também, que, ao contrário do que se esperava, não apareceram abreviaturas próprias do “ internetês ”, relacionado ao uso de redes sociais e emails. b ) aspectos textuais : esta competência é a que se revela como a mais problemática entre os participantes, porque são muitos os problemas observados, desvios acumulados durante toda a formação do estudante : sequência justaposta de ideias sem encaixamentos sintáticos ; redução drástica de estruturas subordinadas, ao lado do aumento na frequência de estruturas coordenadas e absolutas ; redução no uso de conectores para expressar relações lógicas essenciais à construção do texto, substituídas pela exigência de inferência por parte do interlocutor para suprir a sua ausência ; emprego equivocado de operadores que não estabelecem relações lógicas coerentes entre ideias do texto ; emprego inadequado do pronome relativo ( com omissão da preposição ou a utilização de pronome inadequado, como “ onde ” ) ; repetição exaustiva de termos sem a utilização de procedimentos mais sofisticados de substituição ( hiperonímias, hiponímias, nominalizações, expressões metafóricas ) ; frases fragmentadas que comprometem a estrutura lógico-gramatical ; frases formadas apenas por oração subordinada, sem oração principal.
Um importante aspecto a destacar é o baixíssimo desempenho de uma parte dos participantes em relação à estrutura formal do texto produzido, o que é extremamente preocupante ao se levar em conta que são graduandos em fase final de formação.
São frequentes os casos de desvios de estruturação frasal, com uso inadequado ou ausência de conectivos entre parágrafos e entre frases.
Em uma parte dos textos, falta um mínimo de textualidade e de domínio do registro padrão da língua.
Na verdade, observam -se relações linguísticas quase agramaticais, como as estabelecidas pela sequência de gerúndios sem o apoio de um ponto de partida para a organização das informações gramaticais e semânticas.
Observou -se que uma grande parte dos participantes não organizou as ideias em parágrafos, talvez devido ao pequeno número de linhas disponibilizadas para a resposta da questão ou, quem sabe, pela suposição de que não seria necessária essa divisão por não se tratar de um texto no modelo de uma redação dissertativo-argumentativa, como solicitado nos vestibulares.
Quanto à utilização dos mecanismos de referenciação, deve -se destacar a ocorrência, em uma boa parte dos textos, de repetições de palavras ou expressões sem a utilização de termos sinônimos ou pronomes, como seria adequado.
Outro aspecto relevante a ser destacado é a quase total ausência de operadores argumentativos, tanto intrafrasais como interfrasais, repercutindo uma tendência atual da mídia escrita.
Quanto à utilização dos sinais de pontuação, observou -se uma grande precariedade nos textos analisados.
É muito frequente a ocorrência de parágrafos sem marca interna de pontuação para separar os períodos.
Vale observar que não foi penalizada a ausência de vírgula para destacar locuções ou adjuntos adverbiais de pequena extensão deslocados de posição na frase, por ser um uso opcional.
São os seguintes os tipos de problemas encontrados : I. vírgula : utilização de vírgula para separar o sujeito e o predicado ; ocorrência de apenas uma das vírgulas para separar uma palavra, uma expressão ou uma oração encaixada ; uso de vírgula no lugar do ponto para separar ideias que constituem períodos distintos ; II. ponto e vírgula : utilização do ponto e vírgula no lugar de vírgula ; III. ponto final : ausência de ponto final para separar períodos. c ) aspectos morfossintáticos e vocabulares : os resultados são muito transparentes em relação aos aspectos mais problemáticos no desempenho dos participantes.
O desvio mais frequente, em relação à regência, é a falta do sinal indicativo da crase – isso revela que o usuário não tem consciência de que, sob a forma do termo “ a ”, existe a presença de uma preposição “ a ”, exigida pela regência do termo anterior, e de um artigo definido.
Embora em outros exames, como o Enem, a falta de crase seja penalizada em acentuação, nesta avaliação esse desvio foi considerado no âmbito dos aspectos morfossintáticos.
Outro problema relacionado à regência verbal e à nominal, encontrado frequentemente nas questões, foi a ausência de preposição antes de pronome relativo, processo generalizado na modalidade oral da língua, em situações de registro informal.
Apesar da possibilidade de que essa alteração de regência se generalize no padrão escrito da Língua Portuguesa, como já está ocorrendo até em textos jornalísticos, o não emprego da preposição foi penalizado neste processo de avaliação.
Outro desvio muito frequente diz respeito aos processos de concordância verbal e de concordância nominal.
Quanto à concordância de número, observou -se ausência de marca com sujeito posposto ou uso indevido ( uso inadequado da marca de plural comandado pelo núcleo plural da locução adjetiva, apesar de o substantivo que funciona como núcleo do sintagma nominal estar no singular ).
Uma ocorrência generalizada foi a ausência de acento circunflexo na forma plural do presente do indicativo do verbo “ ter ”, que foi considerada como um desvio na concordância verbal e não na acentuação gráfica.
Quanto à concordância de gênero, vários casos foram observados, normalmente no âmbito de sintagmas nominais longos, em que o adjetivo está afastado do substantivo.
Quanto à questão da colocação pronominal, foram poucos os desvios observados.
Apesar de serem aspectos relacionados à oralidade, concluiu -se que, no registro escrito formal, a maioria dos participantes já incorporou regras como a não introdução da frase por um pronome oblíquo e a próclise na presença de um termo atrator.
Não se adotou, entretanto, o padrão excessivamente formal descrito pelas gramáticas normativas em relação à posição do pronome oblíquo em locuções verbais, já que esse uso está muito distante da prática cotidiana, até em textos mais formais.
Quanto aos aspectos vocabulares, alguns tipos de inadequação foram observados, principalmente na Questão 2, relacionada ao funk : expressões da oralidade, seleção vocabular incompatível com o contexto, gerando situações de falta de inteligibilidade ; falta de domínio de vocabulário mais abstrato e de maior Ainda em função da temática desenvolvida na Questão 2, foi significativa a presença de palavras de baixa calão, por retratar um universo cultural específico.
Os textos que continham esse tipo de vocabulário foram considerados inadequados em função do registro formal exigido e foram avaliados com grau zero.
Várias marcas de oralidade foram identificadas, embora não em alta frequência : o uso do pronome relativo “ onde ” como relativo universal, falta de artigo definido antes de substantivo, repetição de palavras por falta de vocabulário, reduções como “ tá ” e “ pra ”, expressões informais. 3.3.2 Componente de Conhecimento Específico Na parte da prova relativa às questões discursivas do Componente de Conhecimento Específico ( Tabela 3.17 ), observa -se que a Média foi mais baixa do que para as questões discursivas do Componente de Formação Geral.
Enquanto no Componente de Formação Geral a Média para estudantes de Tecnologia em Design de Interiores de todo o Brasil foi 54,1, na parte de Conhecimento Específico a Média foi 33,1.
A maior Média deste componente foi obtida pelos estudantes da região Sul ( 38,0 ), e a menor, pelos da região Nordeste ( 30,7 ).
Quanto à variabilidade das notas, o Desvio padrão de todo o Brasil foi 18,8.
O maior Desvio padrão foi encontrado na região Nordeste ( 20,2 ), e o menor, na região Norte ( 16,6 ).
A maior nota Máxima foi obtida na região Sul ( 90,0 ), enquanto a menor nota Máxima foi encontrada na região Norte ( 63,3 ).
Além disso, a nota Mínima ( 0,0 ) foi obtida por pelo menos um aluno em todas as regiões do Brasil, sem exceção.
A Mediana do Brasil como um todo foi 33,3, a mesma obtida nas regiões Norte e Sudeste, enquanto a maior Mediana e a menor Mediana foram obtidas nas regiões Sul ( 38,3 ) e Nordeste ( 30,0 ), respectivamente.
O Gráfico 3.10 representa a distribuição das notas nas questões discursivas no Componente de Conhecimento Específico.
A moda desta distribuição ocorre no intervalo ( 30 ; 40 ].
A análise de cada uma destas questões será feita a seguir. 3.3.2.1 Análise da Questão Discursiva 3 do Componente de Conhecimento Específico Na questão 3, cujos resultados aferidos encontram -se descritos na Tabela 3.18, a Média dos estudantes de todo o Brasil foi 16,9.
O desempenho dos estudantes de todo o Brasil nesta questão foi mais baixo dentre as três questões discursivas desse componente.
A menor Média nessa questão foi obtida pelos alunos da região Norte ( 12,0 ), enquanto a maior Média foi obtida na região Sul ( 20,0 ).
Quanto à variabilidade das notas, o Desvio padrão de todo o Brasil foi 20,7.
O maior Desvio padrão foi obtido na região Nordeste ( 22,3 ), enquanto o menor foi obtido na região Sudeste ( 18,8 ).
A nota Máxima, 95,0 pontos, foi alcançada por, pelo menos, um aluno na região Nordeste, enquanto a menor nota Máxima foi obtida na região Centro-Oeste ( 65,0 ).
A Mediana do Brasil como um todo foi 0,0, o que significa que mais de 50 \% dos alunos obteve nota zero ou deixou a questão em branco, e é a mesma nas regiões Norte e Sudeste.
A maior Mediana nesta questão foi na região Sul ( 15,0 ), e nas regiões Nordeste e Centro-Oeste a Mediana foi 10,0.
A nota Mínima ( 0,0 ) foi obtida por, pelo menos, um estudante em todas as regiões do Brasil, sem exceção.
O Gráfico 3.11 mostra a distribuição das notas na questão discursiva 3, do Componente de Conhecimento Específico da área de Tecnologia em Design de Interiores.
Essa distribuição tem moda na nota zero, com frequência superior a 30 \% do total de alunos.
Destaca -se, também, o intervalo ( 20, 30 ], um máximo local. 3.3.2.2 Comentários sobre as respostas à Questão Discursiva 3 A questão está relacionada diretamente a uma das competências do tecnólogo em Design de Interiores - a elaboração de projetos de interiores - que envolve, no caso, o equacionamento de soluções técnicas, o projeto de um home theater.
Espera -se do estudante a competência de compreender o problema de projeto e orientar quanto à sua solução, reconhecendo a necessidade de buscar a colaboração de especialistas sempre que couber.
É certo, em acréscimo, que qualquer consulta a especialista envolveria, necessariamente, a compreensão do problema, de modo a orientar o encaminhamento de sua solução.
O enunciado da questão, claro e objetivo, solicita a indicação de três projetos complementares à tipologia de um home theater, além da explicitação das características técnicas de cada um deles e da sua contribuição para o desempenho espacial.
Trata -se, portanto, de questão laboriosa, que envolve a compreensão plena da tipologia do ambiente em tela, assim como dos problemas comuns à especificidade de sua função que, notadamente, envolvem o conforto em todos os seus aspectos, a segurança e o atendimento às diretrizes ligadas ao desenvolvimento sustentável.
Observou -se, apesar do enunciado correto, recorrente interpretação equivocada por parte dos estudantes : presente uma imagem meramente ilustrativa de um home theater, muitos deles entenderam terem sido solicitados a apresentar propostas de intervenção naquele espaço específico.
Esta interpretação, no entanto, carece de lógica já que a imagem por si não seria suficiente para fornecer diretrizes para eventuais intervenções, diante do fato de que muitas delas não são aparentes ou não podem ser verificadas desse modo.
Outras recorrentes interpretações equivocadas confundiam o home theater, enquanto tipologia de espaço com função definida, com a aparelhagem de mesma denominação disponibilizada no mercado.
Houve, ainda, inúmeras referências à indicação de outros ambientes nos quais fosse possível instalar o equipamento em tela e, também, a três tipos de desenhos necessários a um projeto de home theater.
Tempo e espaço destinados à resposta foram absolutamente adequados, desde que houvesse competente compreensão da proposta da questão.
Os resultados evidenciam ter sido esta a questão mais difícil entre as discursivas : a média mostrou -se bastante baixa, e o número de questões com pontuação zero, muito alto, não havendo nenhum caso de nota máxima.
O padrão de resposta foi bem definido e adequado à área de Tecnologia do Design de Interiores.
Não houve, no entanto, uma única resposta que atendesse totalmente à sua completude.
Um alto número de respostas totalmente incorretas deveu -se à equivocada interpretação do enunciado, mas mesmo as respostas com notas medianas os estudantes não conseguiram alcançar os nove itens a serem explicitados ( três projetos, principais características e contribuições de cada um ), ou seja, estavam incompletas.
Observou -se, ainda, baixo entendimento quanto ao significado de um projeto complementar, pois, por inúmeras vezes, os estudantes em suas respostas referiam -se a espaços onde se pudesse realizar um projeto de home theater.
Também, foi verificada a confusão entre projeto luminotécnico e projeto elétrico, assim como de conforto acústico e de conforto térmico.
A indicação mais recorrente foi a do projeto luminotécnico, mas as respostas sobre suas características e contribuições foram deficientes, observando -se ausência da terminologia técnica adequada.
O projeto de conforto acústico, também entre os mais citados, muitas vezes reduziu -se a isolamento acústico.
Mesmo assim, foi rara a correta Grande parte das vezes, houve excessiva ênfase no uso de terminologia pouco técnica, como tornar o ambiente “ mais aconchegante ”, “ mais belo ”, “ mais harmônico ” ou “ mais agradável ”, com foco nos aspectos decorativos.
Outra ausência foi a menção à eficiência energética como uma das contribuições de alguns projetos.
Foram muito raras as respostas que a ela se referiam, embora o assunto seja obrigatório nos conteúdos ligados aos projetos de interiores.
O que se observou, de modo repetitivo, foram respostas em que aspectos estéticos foram o foco da resposta, o que aponta para certa prevalência dos mesmos em detrimento dos aspectos técnicos.
Há, sem dúvida, deficiência na formação dos alunos em relação ao conteúdo exigido.
Um profissional de Design de Interiores deve necessariamente conhecer os limites de sua atuação e ter pleno conhecimento dos objetivos a serem alcançados, quando se trata de recorrer a outros profissionais para complementar um projeto.
Cabe, ainda, saber identificar o que solicitar e o que esperar de cada projeto complementar, quando for o caso de buscar expertises correlatas à sua formação.
Esta questão discursiva, que envolve resposta laboriosa, parece encontrar, em grande parte dos concluintes, atitude pouco propícia à sua solução, da mesma forma que as respostas que exigem o uso de terminologia técnica precisa são vistas como difíceis - habituados que estão os estudantes a fazer uso de termos mais genéricos e nada objetivos como “ aconchegante ”, “ agradável ” e “ elegante ”.
Tal fato parece indicar a necessidade de se dar maior \textbf{peso} aos aspectos técnicos e ao uso de terminologia específica compatível com a formação acadêmica desses estudantes.
Deve -se evitar ênfase em aspectos estéticos que, embora igualmente relevantes, em grande parte das vezes atingem uma obviedade que deve, necessariamente, ser evitada na formação profissional. 3.3.2.3 Análise da Questão Discursiva 4 do Componente de Conhecimento Específico A Tabela 3.19 contém as informações relativas à questão 4 do conjunto de questões do Componente de Conhecimento Específico.
A Média geral do Brasil foi 34,2, sendo a menor Média registrada na região Nordeste ( 30,0 ), e a maior, na região Sul ( 40,3 ).
A nota Máxima ( 100,0 ) foi atingida por, pelo menos, um concluinte nas regiões Nordeste e Sul.
Nas regiões Norte e Centro-Oeste a nota Máxima foi 75,0, e na região Sudeste foi 95,0.
A Mediana ( 35,0 ) e a nota Mínima ( 0,0 ) foram as mesmas em todas as regiões, exceto a Mediana na região Sul ( 45,0 ).
O Gráfico 3.12 representa a distribuição de notas da questão discursiva 4, do Componente de Conhecimento Específico.
Essa distribuição apresenta moda no intervalo ( 40 ; 50 ]. 3.3.2.4 Comentários sobre as respostas à Questão Discursiva 4 A questão tem enunciado claro, e a imagem inserida mostra boa visibilidade, sendo pertinente para orientar a resposta pretendida.
Os conhecimentos exigidos e o nível de profundidade são compatíveis com os conteúdos das Diretrizes Curriculares para a Área de tecnologia em Design de Interiores, e abarcam conhecimento compatível com o profissional responsável, tanto pelo projeto, quanto por sua execução, ressalvados os limites de sua atuação.
Sendo relevante a prática em curso dessa natureza, a ocorrência de dupla interpretação não se justifica pois, a linguagem utilizada é precisa.
No entanto, não foi bem compreendida, observando -se como prováveis causas, por parte dos estudantes, a dificuldade de leitura crítica e de interpretação, além de certa desatenção ao enunciado.
Em muitos casos, tais fatos foram responsáveis pela ausência de pontuação.
Verificou -se a ocorrência de muitas respostas equivocadas que se referiam a etapas de projeto e não a etapas de execução, interpretação totalmente injustificada.
O tempo de resposta e o espaço destinado à mesma mostraram -se adequados.
A questão pode ser considerada fácil e dentre as três discursivas posiciona -se em termos medianos.
Os estudantes que tiveram contato com a prática, provavelmente foram os mais bem-sucedidos.
O padrão de resposta pode ser identificado nos textos dos estudantes, embora o nível de detalhamento das etapas se apresentasse, na maior parte das vezes, mais resumido, algumas delas sendo suprimidas ou subentendidas em respostas mais genéricas.
Não ocorreram respostas não previstas.
Mencione -se certa recorrência, cerca de 10 \%, de respostas equivocadas que descreviam as etapas do projeto de interiores, em detrimento das etapas do projeto executivo pertinentes àquele caso específico.
E, ainda, uma clara e correta sequência de um projeto executivo, porém apartado do caso apresentado, onde se inseriam, inclusive, itens não constantes do mesmo.
Como o enunciado não se referia ao revestimento usado nas paredes, pouquíssimos estudantes citaram a etapa da pintura, embora a mesma pudesse ser inferida.
Outros mencionaram as etapas de uma obra aleatória, incluindo itens que o enunciado deixava claro que não serviam para aquele caso, como a construção de paredes, a colocação do piso, além de instalações hidrossanitárias.
Na correção, quando citadas, tais etapas foram desconsideradas.
A maior parte dos candidatos citou entre 3 e 4 etapas ( dentre as 8 exigidas ), na sequência correta.
As etapas recorrentemente ignoradas foram as de proteção do piso, de demolição, de pintura e de limpeza ( etapas 1, 2, 5 e 8 conforme o padrão de resposta ).
Com isso, a maior parte das notas foi mediana.
Alguns erros parecem inaceitáveis para designers prestes a se tornarem profissionais, como por exemplo, a utilização do termo “ imobiliário ” ao invés de “ mobiliário ”, ou erros ortográficos em estrangeirismos - embora termos de uso corrente na Área, como layout e briefing.
Nesta questão, o desempenho dos estudantes mostrou -se insuficiente em relação ao seu baixo nível de complexidade e por referir -se a conteúdo cerne da atuação profissional.
Foram raríssimos os candidatos que citaram as oito etapas previstas no padrão de resposta e, ainda mais raros os que alcançaram a pontuação máxima, ao também descreverem a sequência correta das etapas.
A resposta mais recorrente, acarretando nota zero, foi a citação de etapas de projeto e não etapas de execução da obra, interpretação inaceitável.
A maioria dos estudantes citou entre 3 ou 4 etapas de execução, em ordem muitas vezes correta, consideradas e pontuadas mesmo quando intercaladas de etapas equivocadas.
Não houve diversidade de soluções apresentadas, exceto as que pautadas em interpretações inadequadas.
Para um formando, o conhecimento das etapas de uma obra e sua sequência de execução constitui informação imprescindível.
É possível que o resultado mediano alcançado evidencie uma formação não suficientemente conectada com a prática de obra, o que, no caso de Cursos Superiores de Tecnologia, exige particular atenção.
Os textos corridos dificultaram a correção : muitas vezes o estudante listava os itens em sequência e, ao final acrescentava : ” mas antes disso tudo... ”, inserindo um item anterior aos demais.
A questão, de padrão fácil, só não obteve resultados melhores por abarcar aspecto prático da profissão nem sempre vivenciado na formação acadêmica.
O resultado aponta para que se busque, nos meios acadêmicos, um olhar mais direcionado ao mundo da execução, pois se trata de tarefa de muita relevância na práxis profissional. 3.3.2.5 Análise da Questão Discursiva 5 do Componente de Conhecimento Específico A Tabela 3.20 contém as informações relativas à questão 5 do conjunto do Componente de Conhecimento Específico.
O desempenho dos estudantes nessa questão foi superior ao das questões 3 e 4.
A nota Média dos estudantes de todo o Brasil foi 48,1.
A maior Média foi registrada na região Sul ( 53,8 ), enquanto a menor Média foi registrada na região Nordeste ( 43,2 ).
Quanto à variabilidade das notas, o Desvio padrão dos alunos do Brasil, como um todo, foi 34,2.
Enquanto o maior desvio foi encontrado na região Nordeste ( 35,3 ), o menor foi encontrado na região Norte ( 32,8 ).
A nota Máxima ( 100,0 ) e a nota Mínima ( 0,0 ) são as mesmas em todas as regiões, exceto a nota Máxima na região Norte ( 90,0 ).
Para o conjunto de alunos de Tecnologia em Design de Interiores do Brasil, a Mediana foi 55,0, sendo a mesma obtida na região Sudeste, enquanto nas demais regiões a Mediana foi : 60,0 nas regiões Norte e Sul ; 45,0 na região Nordeste ; e 40,0 na região Centro-Oeste.
O Gráfico 3.13 representa a distribuição das notas da questão discursiva 5 do Componente de Conhecimento Específico.
Essa distribuição tem moda no intervalo ( 80 ; 90 ], apresentando os intervalos da nota zero, ( 20, 30 ] e ( 50, 60 ] como máximos locais. 3.3.2.6 Comentários sobre as respostas à Questão Discursiva 5 Trata -se da questão, certamente, de menor complexidade entre as discursivas.
O enunciado é claro e o nível de profundidade adequado às Diretrizes Curriculares da Área de Tecnologia em Design de Interiores : aborda questão de cunho teórico ligada ao ritmo na solução de composição espacial, tendo conteúdo pertinente a um futuro profissional apto a projetar ambientes interiores.
Em relação ao tempo de resposta e ao espaço disponível não foi observada nenhuma incongruência.
A média alcançada pelos estudantes, mais alta entre as três discursivas, ratifica ser esta a questão de maior facilidade entre as discursivas.
Foi também a questão em que se obteve ainda o maior número de respostas com pontuação máxima.
O padrão de respostas, principalmente no que se refere à figura 1, foi restrito, disponibilizando apenas três respostas certas para os três itens requeridos.
Para a figura 2, eram quatro as possibilidades de respostas corretas, sendo que o estudante deveria fornecer três delas.
Muitas vezes as respostas eram imprecisas, por exemplo, referindo- se simplesmente ao “ piso ou ao teto ”, em preferência à “ repetição / composição das linhas no piso e no teto ”.
Não se atribuiu pontuação nesses casos, pois fora solicitado que indicassem características que representassem o ritmo e não o local de aplicação das mesmas.
Foi recorrente a descrição de algumas características não contempladas no padrão de respostas, sobretudo nas respostas consideradas medianas.
Para a figura 1, foram citadas com frequência : linhas retas do ambiente, padrões geométricos das portas, repetição da forma circular na luminária, repetição das linhas dos degraus, livros alinhados na estante.
Para a figura 2, foram citadas com frequência : iluminação entre as linhas do teto, ondulações na parede do fundo do ambiente, repetição das louças sobre a mesa, repetição das linhas curvas das cadeiras, linhas retas dos mobiliários em contraposição às sinuosas no teto.
Alguns equívocos evidenciaram desatenção em atender ao total de três características ( em cada uma das figuras ) e em alguns casos as respostas denunciavam claramente a incompreensão quanto ao elemento ritmo.
Porém, parte significativa dos erros cometidos estavam relacionados diretamente à dificuldade de leitura e interpretação do enunciado além de deficiência na expressão escrita, pois os estudantes foram frequentemente sucintos, além do necessário para atender ao solicitado.
Por vezes, ainda, ambíguos em suas respostas : recorrentemente inferiu -se confusão entre os sentidos dos termos ritmo, harmonia, estilo e unidade, sendo o termo harmonia, presente no enunciado, muitas vezes usado como sinônimo de ritmo.
Ocorreram ainda respostas que descreviam o estilo dos ambientes das imagens, através de uma terminologia como art déco, moderno, clean, dentre outros.
Constatou -se que parte dos estudantes evidenciava conhecer o significado de ritmo no ambiente - o principal objetivo da investigação da questão - mas em muitos casos as pontuações foram aquém do seu conhecimento a este respeito, por não terem descrito de forma clara e precisa os itens previstos no padrão de respostas.
O baixo resultado das respostas consideradas fracas ocorreu, desse modo, por deficiência descritiva e ausência de clareza na expressão da característica.
Quando evidente a falta de compreensão sobre o significado do ritmo na composição espacial, caso de pouco menos da metade dessas notas, a pontuação aplicada foi zero.
Referência às “ linhas de teto e piso ” foram as respostas mais recorrentes para a figura 1, enquanto as “ ondulações de teto e colunas ” igualmente o foram para a figura 2.
Os principais erros abarcaram a dificuldade em atender ao comando, identificando “ características que representassem o ritmo ”.
Quanto às respostas que descreviam os ambientes, ora indicando estilos, como vintage, ora suas características como “ confortável ”, “ elegante ”, “ descontraído ”, “ formal ”, entre outras, essas não foram pontuadas.
Observe -se aqui a recorrente presença de uma leitura baseada em \textbf{observações} por vezes excessivamente triviais, que permearam muitas respostas.
Em parte, verificou -se que deficiência interpretativa, dificuldade de expressão escrita e falta de objetividade foram fatores que prejudicaram severamente o desempenho nessa questão.
Os resultados evidenciam a necessidade de se fortalecer, na formação profissionais da Área de Design de Interiores, a competência para interpretar textos escritos e a capacidade de expressão, usando termos adequados à formação profissional, abandonando certos vícios descritivos pautados em análises banalizadas dos espaços interiores. 3.3.3 Considerações Finais A participação da Área de Tecnologia em Design de Interiores no Enade/2015 permite que os resultados das três questões discursivas ofereçam um significativo olhar sobre a formação acadêmica, em campo de saber ainda pouco delimitado em relação a outras Áreas da Educação Superior.
É possível apresentar alguns pontos onde a cuidadosa intervenção dos dirigentes e docentes das Instituições de Ensino possa provocar mudança de qualidade, não somente capaz de melhorar a formação acadêmica, como contribuir sobremaneira para o processo de consolidação da Área de Design de Interiores.
O uso e a assimilação de uma linguagem técnica pertinente às adequadas intervenções nos espaços interiores aptas a promover saúde, segurança e qualidade de vida aos usuários podem constituir um passo fundamental.
Ao designer, cabe apontar soluções de naturezas diversas que certamente envolvem aspectos sintáticos, semânticos e pragmáticos do uso dos ambientes, e que não podem, sob nenhum ponto de vista, restringir -se ao discurso mercadológico e banalizado, comumente observado na mídia não especializada, mas observado com muita frequência nas respostas dos estudantes.
Outro aspecto a considerar envolve a ênfase na leitura e interpretação de textos o que ampliaria a competência da linguagem escrita, ponto que contribuiria substancialmente para o desenvolvimento do profissional que, além do projeto gráfico, faz uso de peças escritas que o complementam - como memoriais descritivos, justificativos e especificações técnicas, dentre outras.
Mesmo que o fazer em curso de tecnologia seja privilegiado, não há por que se admitir a manutenção do reduzido hábito de leitura oriundo do Ensino Médio, como se sabe acontecer.
Uma terceira questão aponta ainda para a prática profissional no que diz respeito ao projeto executivo e à execução da obra, tarefa das mais árduas para as instituições de ensino, pois os estágios nem sempre suprem esse conhecimento, grande parte das vezes mantendo os estagiários apartados do mundo da obra e focados na tela de computadores em programas de desenho.
O papel da supervisão docente nesses programas de estágio precisa buscar suprir esse viés, entendendo a prática como instrumento importante na formação do estudante.
Tais considerações, oriundas da experiência na correção dessas questões discursivas em prova de desempenho de estudantes concluintes, evidenciam o seu \textbf{peso} como elemento que pode orientar as diversas os cursos rumo a um padrão de excelência.
No entanto, é preciso que os próprios estudantes compreendam e creditem esse papel ao Enade/2015, pois a deficiência de conhecimento evidenciada em muitas respostas pode, na verdade, também indicar o desvalor com que a prova é realizada.
Faz -se necessário, portanto, refletir sobre os mecanismos que poderiam contribuir para sua mudança de significado junto ao público alvo.

% matched lemmas: desordem, exacerbar, fraternidade, intérprete, liberal, morfema, nocivo, novembro, observação, peso
\end{textsample}
